\documentclass[11pt]{article}
\usepackage[a4paper,margin=1in]{geometry}
\usepackage{amsmath,amssymb,mathtools,bm}
\usepackage{enumitem,booktabs,array}
\usepackage{tcolorbox}
\usepackage{hyperref}
\usepackage[protrusion=true,expansion=false]{microtype}
\hypersetup{colorlinks=true,linkcolor=blue,urlcolor=blue,citecolor=blue}
\setlist[itemize]{topsep=2pt,itemsep=2pt}
\setlist[enumerate]{topsep=2pt,itemsep=2pt}
\newtcolorbox{keybox}{colback=blue!3!white,colframe=blue!40!black,boxrule=0.4pt,arc=1mm}
\newtcolorbox{alertbox}{colback=red!3!white,colframe=red!40!black,boxrule=0.4pt,arc=1mm}
\newtcolorbox{answerbox}{colback=green!3!white,colframe=green!40!black,boxrule=0.4pt,arc=1mm}
\newtcolorbox{drillbox}{colback=yellow!5!white,colframe=orange!60!black,boxrule=0.4pt,arc=1mm}
\newcommand{\EV}{\mathbb{E}}
\newcommand{\Var}{\mathrm{Var}}
\newcommand{\Cov}{\mathrm{Cov}}
\newcommand{\blank}[1]{\underline{\hspace{#1}}}

\title{E01-02: Klein, Crawford \& Alchian (1978, JLE) \\
\large ``Vertical Integration, Appropriate Rents, and the Competitive Contracting Process'' --- Active Recall Workbook}
\author{Empirical Corporate Finance II --- Exam Prep}
\date{\today}
\begin{document}\maketitle

\begin{keybox}
\textbf{Paper in one sentence.} When an asset is ``relationship-specific'' --- worth more inside a particular trading relationship than outside --- the gap (the \emph{quasi-rent}) creates a holdup temptation for the partner, so parties choose vertical integration or long-term contracts to protect the quasi-rent from opportunistic appropriation.

\textbf{Placement.} Concretizes Coase (1937): names the transaction cost that matters (holdup of relationship-specific investment). Direct antecedent to Williamson (1979, 1985) on asset specificity, Grossman-Hart (1986) on residual rights, and modern empirical make-or-buy work.
\end{keybox}

\section*{Round 1: Complete Walk-Through}

\subsection*{1.1 Key concepts}
\begin{itemize}
\item \textbf{Specific asset.} An asset with a lower value in its second-best use than in its current one.
\item \textbf{Appropriable quasi-rent.} (Value in current use) $-$ (value in best alternative use). This is the rent a partner can attempt to capture via opportunistic renegotiation.
\[
\text{Quasi-rent} = V_{\text{current}} - V_{\text{alt}} \ge 0.
\]
\item \textbf{Holdup.} Post-investment, the non-investing partner renegotiates to capture the quasi-rent, leaving the investor with only the alternative value.
\end{itemize}

\subsection*{1.2 Canonical example: Fisher Body / General Motors}
Fisher Body invests in auto-body stamping dies specific to GM models. Once sunk, Fisher cannot redeploy the dies profitably elsewhere. GM can threaten to reduce the price it pays and capture Fisher's quasi-rents. Result: GM vertically integrates Fisher Body in 1926 (acquires ownership of the specific assets and the relationship).

\subsection*{1.3 Governance choice}
To protect quasi-rents, parties choose among:
\begin{itemize}
\item Spot-market transactions (only when specificity is low).
\item Long-term contracts (when specificity exists but contract is writable).
\item Vertical integration (when specificity is high and contracts are incomplete).
\end{itemize}

\subsection*{1.4 Formal structure (intuitive)}
Let $k$ denote relationship-specific investment, $V(k)$ current-relationship value, $\bar V$ alternative-use value. Ex ante surplus is $V(k) - k$. After $k$ is sunk, ex post bargaining splits $V(k) - \bar V$. If the investor fears holdup, $k^*$ is underinvested relative to first-best. Integration/long-term contract restores incentives to invest.

\subsection*{1.5 Predictions}
\begin{itemize}
\item Higher asset specificity $\Rightarrow$ more vertical integration / long-term contracting.
\item Contracts substitute for integration when contingencies are foreseeable and enforceable.
\item Spot markets dominate for generic (non-specific) inputs.
\end{itemize}

\subsection*{1.6 Influence}
\begin{itemize}
\item Williamson: asset specificity becomes the central dimension of transaction-cost economics.
\item Joskow (1987) on coal contracts, Monteverde-Teece (1982) on auto parts empirically support the predictions.
\item Grossman-Hart: shifts focus to \emph{who owns} residual control rights rather than whether a contract exists.
\item Corporate-finance integration/outsourcing decisions; franchising vs.\ company-owned.
\end{itemize}

\subsection*{1.7 Caveats}
\begin{alertbox}
\begin{itemize}
\item Coase-Fisher-Body narrative has been challenged historically (e.g., Casadesus-Masanell-Spulber 2000 dispute the holdup account).
\item Theory predicts integration but cannot precisely pin down boundary --- contract vs.\ integration margin is empirical.
\item Assumes opportunism; repeated-game cooperation may substitute for integration in practice.
\end{itemize}
\end{alertbox}

\section*{Round 2: Level 1 Fill-in}
\begin{enumerate}
\item Appropriable quasi-rent = value in current use $-$ value in \blank{2cm} use.
\item Holdup arises \emph{after} relationship-specific \blank{1.5cm} is sunk.
\item Canonical example: \blank{2cm} Body and General Motors.
\item Three governance choices: spot market, long-term \blank{1.5cm}, vertical \blank{2cm}.
\item Key later author who built on KCA: \blank{2cm} (1979, 1985).
\end{enumerate}

\section*{Round 3: Level 2 Prompts}
\begin{enumerate}
\item Define appropriable quasi-rent and asset specificity; work through the Fisher-GM example.
\item Explain the holdup problem and why relationship-specific investment is underprovided.
\item Derive the governance-choice prediction: specificity $\Rightarrow$ integration.
\item Relate KCA to Coase (1937) and to Williamson (1979).
\item Discuss empirical tests (Joskow, Monteverde-Teece) and limitations.
\end{enumerate}

\section*{Round 4: Minimal Scaffolding}
From a blank page: (i) quasi-rent definition, (ii) Fisher-GM example, (iii) holdup logic, (iv) governance prediction, (v) descendants.

\section*{Round 5: Blank Page}
Essay: relationship-specific investment and the theory of vertical integration, KCA's contribution.

\vspace{6pt}
\begin{drillbox}
\textbf{Speed Drill.} (1) Quasi-rent formula. (2) Canonical example. (3) Governance continuum. (4) Holdup problem statement. (5) One empirical test.
\end{drillbox}

\section*{Quick Reference \& Answer Key}
\begin{answerbox}
\begin{itemize}
\item \textbf{Quasi-rent.} $V_{\text{current}} - V_{\text{alt}}$.
\item \textbf{Example.} Fisher Body / GM 1926.
\item \textbf{Logic.} Specificity $\Rightarrow$ holdup risk $\Rightarrow$ integration.
\item \textbf{Continuum.} Spot $<$ long-term contract $<$ integration.
\end{itemize}
\end{answerbox}

\begin{keybox}
\textbf{30-second exam pitch.} Klein, Crawford \& Alchian (1978) give transaction-cost theory its sharpest mechanism: when an asset is relationship-specific --- worth more inside a particular trading relationship than outside --- there is an \emph{appropriable quasi-rent} equal to the gap between current-use and alternative-use value. Once the investment is sunk, the non-investing partner can renegotiate and capture that quasi-rent --- the holdup problem. Anticipating this, parties choose governance to protect the rent: spot markets for generic assets, long-term contracts when contingencies are writable, and vertical integration when specificity is high and contracts incomplete. The canonical story is Fisher Body / GM (1926), where GM integrated Fisher to internalize specific stamping-die investments. The paper names the transaction cost that matters, launches Williamson's asset-specificity program, and underlies modern empirical make-or-buy and outsourcing research. In corporate finance, it motivates why firms own rather than rent critical specific assets.
\end{keybox}

\end{document}
